%=====================================================================
%  VIRTUAL SYSTEMS AND SERVICES --- FALL 2026
%  One-page weekly lecture schedule
%
%  SOURCE OF RECORD: parts/G_schedule.tex in the handout.  Change the
%  handout first, then this page -- the slide pipeline parses the
%  appendix, not this file.  Content modernises HEC 2017
%  (Curriculum of BS & MS CS, SE & IT, pp. 71-130): same course, same
%  credit hours and prerequisite, hypervisors carried forward into
%  containers, Kubernetes and cloud-native virtualization.  See
%  ../HEC AI-DS Curriculum Map/ for the edition-by-edition trace (2017
%  -> unspecified in 2023 -> dropped in 2025, replaced by Cloud
%  Computing and Microservices Architecture and Docker Containers).
%
%  Build:  pdflatex VSS_Fall2026_Lecture_Schedule.tex
%=====================================================================
\documentclass[10pt,a4paper]{article}
\usepackage[a4paper,left=1.2cm,right=1.2cm,top=0.85cm,bottom=0.6cm]{geometry}
\usepackage[T1]{fontenc}
\usepackage{lmodern}
\usepackage{microtype}
\usepackage[table]{xcolor}
\usepackage{array,ragged2e}
\usepackage{graphicx}
\usepackage{fontawesome5}
\pagestyle{empty}
\pdfinfo{/Title (Virtual Systems and Services, Fall 2026: Weekly Lecture Schedule)}
\setlength{\parindent}{0pt}
\setlength{\parskip}{0pt}

%--- palette: shared BSIT palette (../Data Science Fall 2026/dshandout.sty
%    family); part colours darkened a little so white chip text stays
%    legible at 6 pt
\definecolor{primary}{HTML}{1A5276}
\definecolor{secondary}{HTML}{2E86C1}
\definecolor{accent}{HTML}{C0392B}
\definecolor{lightbg}{HTML}{EBF5FB}
\definecolor{neutral}{HTML}{707B7C}
\definecolor{sub}{HTML}{465152}
\definecolor{p1}{HTML}{1A5276}
\definecolor{p2}{HTML}{2874A6}
\definecolor{p3}{HTML}{1E8449}
\definecolor{p4}{HTML}{8E44AD}
\definecolor{p5}{HTML}{117A65}
\definecolor{p6}{HTML}{8A6D0A}

%--- building blocks -------------------------------------------------
\newcommand{\chip}[2]{{\setlength{\fboxsep}{1.3pt}\colorbox{#1}{\makebox[1.45em][c]{%
  \fontsize{6.4}{7}\selectfont\bfseries\color{white}#2}}}}
\newcommand{\chipw}[2]{{\setlength{\fboxsep}{1.3pt}\colorbox{#1}{%
  \fontsize{6.4}{7}\selectfont\bfseries\color{white}\,#2\,}}}
% \ch{part}{number}{title}{what it covers}  -- one chapter, own paragraph
\newcommand{\ch}[4]{\chip{p#1}{#2}\,\textbf{#3}\ {\color{neutral}--}\ {\fontsize{7.3}{8.5}\selectfont\color{sub}#4}\par}
\newcommand{\ex}[1]{{\bfseries\color{accent}#1}}      % exams and the big deliverables
\newcommand{\ev}[1]{{\bfseries\color{primary}#1}}     % quizzes, assignments, project points
\newcommand{\none}{{\color{neutral}---}}
\newcommand{\hdr}[1]{{\bfseries\color{white}#1}}
\newcommand{\lab}[1]{{\color{sub}#1}}

\newlength{\topicw}   % set inside the document, once \linewidth and \tabcolsep are final
\newcolumntype{K}{>{\centering\arraybackslash}p{0.75cm}}
\newcolumntype{T}{>{\RaggedRight\arraybackslash}p{\topicw}}
\newcolumntype{L}[1]{>{\RaggedRight\arraybackslash}p{#1}}
% \leavevmode: a cell that opens with \color puts a whatsit before its paragraph, and
% \vtop then aligns it by its top edge -- one strut lower than a cell that opens with text
\newcommand{\wk}[4]{\leavevmode{\large\bfseries\color{primary}#1} & \leavevmode #2 &
  \leavevmode #3 & \leavevmode #4 \\ \hline}

\newcommand{\fact}[3]{{\setlength{\fboxsep}{4pt}\colorbox{lightbg}{%
  \parbox[t]{\dimexpr#1-8pt\relax}{\RaggedRight\fontsize{6.2}{7}\selectfont\bfseries\color{neutral}\MakeUppercase{#2}\\[1.5pt]%
  \fontsize{7.8}{9.2}\selectfont\bfseries\color{primary}#3}}}}
\newcommand{\tile}[3]{{\colorbox{lightbg}{\parbox[c][0.9cm][c]{\dimexpr(\linewidth-3pt*\numexpr#1-1\relax)/#1-6pt-0.5pt\relax}{%
  \centering{\large\bfseries\color{primary}#2}\\[-0.5pt]{\fontsize{6.6}{7.4}\selectfont\color{sub}#3}}}}}

% \legend{I}{II}{III}{IV}{V}{VI} -- the six part names, as a 3 x 2 grid of colour chips
\newcommand{\legend}[6]{{\fontsize{6.6}{8}\selectfont\color{sub}%
  \begin{tabular}{@{}p{0.31\linewidth}p{0.31\linewidth}p{0.31\linewidth}@{}}
  \chipw{p1}{I}\ #1 & \chipw{p2}{II}\ #2 & \chipw{p3}{III}\ #3 \\[1pt]
  \chipw{p4}{IV}\ #4 & \chipw{p5}{V}\ #5 & \chipw{p6}{VI}\ #6
  \end{tabular}}}

\begin{document}
\fontsize{8.2}{9.8}\selectfont

%--- header ----------------------------------------------------------
\begin{minipage}[c]{1.5cm}%
  \IfFileExists{assets/iub-logo.png}{\includegraphics[height=1.15cm]{assets/iub-logo.png}}{}%
\end{minipage}\hfill
\begin{minipage}[c]{0.65\linewidth}\centering
  {\footnotesize\bfseries\color{primary}THE ISLAMIA UNIVERSITY OF BAHAWALPUR}\\[1pt]
  {\scriptsize\color{neutral}Department of Information Technology}
\end{minipage}\hfill
\begin{minipage}[c]{1.5cm}\raggedleft
  \IfFileExists{assets/dit-logo.png}{\includegraphics[height=1.15cm]{assets/dit-logo.png}}{}%
\end{minipage}

\vspace{3pt}
{\color{secondary}\rule{\linewidth}{1.4pt}}\\[2pt]
{\Large\bfseries\color{primary}VIRTUAL SYSTEMS AND SERVICES}\hfill
{\normalsize\bfseries\color{secondary}Weekly Lecture Schedule\ \textbullet\ Fall 2026}\\[1pt]
{\footnotesize\color{sub}Bachelor of Science in Information Technology\ \textbullet\
 \emph{From Hypervisors to the Cloud-Native Stack}}\\[1pt]
{\footnotesize\color{sub}Modernised from the HEC 2017 outline: the same hypervisor foundations, carried forward into containers, Kubernetes and cloud virtualization}\\[1pt]
{\color{secondary}\rule{\linewidth}{0.6pt}}

\vspace{4pt}
\fact{0.235\linewidth}{Credit hours}{3 (2--1): 2 lecture h + 1 lab h a week}\hfill
\fact{0.235\linewidth}{Duration}{16 teaching weeks; midterm Week 8, project due Week 16}\hfill
\fact{0.235\linewidth}{Prerequisite}{Programming Fundamentals (HEC 2017 outline)}\hfill
\fact{0.235\linewidth}{Contents}{23 chapters in 6 parts; hypervisors, containers and cloud VMs}

\vspace{5pt}
%--- the schedule ----------------------------------------------------
\begingroup   % keep \rowcolors and the table geometry from leaking into the legend
\renewcommand{\arraystretch}{1.3}
\setlength{\tabcolsep}{3.5pt}
\setlength{\topicw}{\dimexpr\linewidth-0.75cm-4.05cm-3.3cm-6\tabcolsep\relax}
\arrayrulecolor{primary!18}
\rowcolors{2}{white}{lightbg}
\noindent\begin{tabular}{@{}K T L{4.05cm} L{3.3cm}@{}}
\rowcolor{primary}
\hdr{\faIcon{calendar-alt}} & \hdr{\faIcon{book-open}\ \ Lecture topics} &
\hdr{\faIcon{server}\ \ Laboratory (1 h)} & \hdr{\faIcon{clipboard-check}\ \ Assessment} \\
\wk{1}
 {\ch{1}{1}{What Virtualization Is}{abstraction of compute, storage and network; why virtualize; from IBM VM/370 to the cloud}
  \ch{1}{2}{Virtualization Architectures}{Type-1 vs Type-2 hypervisors, full virtualization, paravirtualization, hardware-assisted virtualization}}
 {\lab{Install a Type-2 hypervisor; a first VM and snapshot}}{\none}
\wk{2}
 {\ch{1}{3}{The Virtual Machine Monitor}{the Popek--Goldberg theorem, trap-and-emulate, binary translation, ring deprivileging}}
 {\lab{Trace a privileged instruction; trap-and-emulate by hand}}{\none}
\wk{3}
 {\ch{2}{4}{Hardware Support for Virtualization}{Intel VT-x/VT-d, AMD-V/AMD-Vi, extended/nested page tables, the IOMMU}}
 {\lab{Check VT-x/AMD-V flags; enable nested virtualization}}{\none}
\wk{4}
 {\ch{2}{5}{Major Hypervisor Platforms}{VMware ESXi/vSphere, Microsoft Hyper-V, KVM/QEMU, Xen; choosing between them}}
 {\lab{The same workload on ESXi, Hyper-V and KVM}}{\none}
\wk{5}
 {\ch{2}{6}{CPU and Memory Virtualization}{vCPU scheduling, overcommit, ballooning, transparent page sharing, NUMA}}
 {\lab{Ballooning and overcommit under memory pressure}}{\none}
\wk{6}
 {\ch{2}{7}{Paravirtualization and Device Virtualization}{virtio and PV drivers, emulated vs passthrough devices, SR-IOV, VFIO}}
 {\lab{virtio-net against an emulated NIC, throughput compared}}{\ev{Quiz 1} (Ch 1--7)}
\wk{7}
 {\ch{3}{8}{From Virtual Machines to Containers}{namespaces, cgroups, union filesystems, OS-level vs hardware virtualization}
  \ch{3}{9}{Docker and the Container Ecosystem}{images and layers, the OCI spec, registries, writing a sound Dockerfile}}
 {\lab{Namespaces and cgroups by hand; a first Dockerfile}}{\ev{Assignment 1 set}}
\wk{8}
 {\ch{3}{10}{Container Orchestration with Kubernetes}{pods, deployments and services, scheduling, self-healing, scaling}}
 {\lab{Midterm, then a Deployment and a Service on kind}}{\ex{Midterm} (Ch 1--9)}
\wk{9}
 {\ch{3}{11}{Lightweight and Secure Virtualization}{microVMs (Firecracker), gVisor, Kata Containers, unikernels}}
 {\lab{A Firecracker microVM, boot time measured}}{\ev{Project set}\newline \ev{Assignment 1 due}}
\wk{10}
 {\ch{4}{12}{Storage Virtualization}{SAN/NAS abstraction, virtual disks, thin provisioning, software-defined storage, Ceph/vSAN}
  \ch{4}{13}{Network Virtualization and SDN}{virtual switches, VLANs and VXLAN overlays, SDN controllers, NFV}}
 {\lab{A thin-provisioned pool; a small Ceph cluster}}{\ev{Project proposal due}}
\wk{11}
 {\ch{4}{14}{Cloud Service and Deployment Models}{IaaS/PaaS/SaaS/FaaS, public/private/hybrid/multi-cloud, VM instances on AWS/Azure/GCP}}
 {\lab{One workload, three clouds: AWS, Azure, GCP instances}}{\ev{Quiz 2} (Ch 10--14)}
\wk{12}
 {\ch{4}{15}{Desktop and Application Virtualization}{VDI, Citrix/Horizon, application streaming, remote browser isolation}
  \ch{5}{16}{Virtualization Management Platforms}{vCenter, Hyper-V Manager/SCVMM, oVirt, Proxmox, OpenStack, Terraform}}
 {\lab{A VDI session; a streamed application}}{\ev{Assignment 2 set}}
\wk{13}
 {\ch{5}{17}{High Availability, Live Migration and DR}{vMotion/Live Migration, clustering, snapshots, checkpointing, backup strategy}}
 {\lab{Live Migration with zero dropped pings}}{\ev{Project experiments frozen}}
\wk{14}
 {\ch{5}{18}{Performance Monitoring and Capacity Planning}{hypervisor metrics, the noisy-neighbour effect, benchmarking virtualized workloads}
  \ch{5}{19}{Security in Virtualized Environments}{VM escape, hypervisor attack surface, Spectre/Meltdown in multi-tenant clouds, container hardening}}
 {\lab{fio/stress-ng: native vs virtualized; a container break-out drill}}{\ev{Quiz 3} (Ch 15--19)\newline \ev{Assignment 2 due}}
\wk{15}
 {\ch{6}{20}{GPU and Accelerator Virtualization}{GPU passthrough, vGPU and MIG partitioning, virtualizing AI/ML workloads}
  \ch{6}{21}{Confidential Computing and Emerging Hardware}{Intel TDX, AMD SEV, trusted execution environments, composable infrastructure}}
 {\lab{A vGPU/MIG partition for a training job}}{\none}
\wk{16}
 {\ch{6}{22}{Case Studies in Enterprise and Cloud Virtualization}{data-centre consolidation, cloud migration, edge virtualization}
  \ch{6}{23}{Semester Project Presentations}{deployed-service demos, peer review, course review}}
 {\lab{Peer walkthroughs of each team's deployed service}}{\ex{Project report and presentation}}
\end{tabular}
\endgroup

\vspace{4pt}
%--- assessment weights ----------------------------------------------
{\fontsize{6.6}{8}\selectfont\bfseries\color{neutral}\MakeUppercase{Assessment weights}}\\[2pt]
\tile{3}{20\%}{Sessional: quizzes, assignments, project}\hfill
\tile{3}{30\%}{Mid-term, Week 8}\hfill
\tile{3}{50\%}{Final-term, all chapters}

\vspace{3pt}
%--- legend and notes ------------------------------------------------
\legend{Foundations of Virtualization}{Hypervisors and Compute Virtualization}{Containers and Lightweight Virtualization}
       {Storage, Network and Cloud Virtualization}{Management, Performance and Security}{Emerging Hardware and Practice}

\vspace{1pt}
{\fontsize{6.3}{7.6}\selectfont\color{neutral}%
Modernised from HEC 2017 (\emph{Curriculum of BS \& MS CS, SE \& IT}, pp.\ 71--130): the original hypervisor trio --- VMware, Xen and Hyper-V --- opens the course in Weeks 3--6; Weeks 7--16 carry the same ground into containers, Kubernetes, cloud models and confidential computing, anticipating HEC 2025's \emph{Cloud Computing} and \emph{Microservices Architecture and Docker Containers}, the two courses that replaced it. The 2017 set texts (Stanney, Burdea) describe virtual reality, not system virtualization, and are not used here. Chapter numbers refer to the course handout, \emph{Virtual Systems and Services --- From Hypervisors to the Cloud-Native Stack} (239 pp., 23 chapters, 7 appendices); Appendix~G of that handout is the source of record for this page.
Weights follow the department's standard split: Sessional 20\%, Mid-term 30\%, Final-term 50\%.\par}

\end{document}
